%% Comparison of LaRT_v2.00 with RT-scat (Kym et al. 2026), and problems found in RT-scat
%% Author: Kwang-Il Seon
\documentclass[english,a4paper]{my_memo}
\synctex=1
\usepackage{amsmath}
\usepackage{amssymb}
\usepackage{booktabs}
\usepackage{graphicx}
\usepackage{microtype}
\usepackage{babel}
\graphicspath{{figures/}}

\begin{document}

\title{LaRT and RT-scat: Clumpy and Rotating Ly$\alpha$ Models of Kym et al.\ (2026)}
\author{Kwang-Il Seon}
\date{\docmoddate}
\maketitle

\tableofcontents
\bigskip

%=============================================================================
\section{Purpose}
%=============================================================================

Kym, Chang, Gronke \& Seon (2026, arXiv:2609.08396) computed Ly$\alpha$ spectra of
rotating smooth and clumpy spheres with the code RT-scat (S.-J.~Chang).  While
recomputing their Figure~9 with LaRT\_v2.00, the two codes were compared directly on
a static clumpy sphere, and the smooth rotating sphere was examined for numerical
artifacts.  This note records the setup, the measured results, the problems found in
RT-scat, and the changes made to LaRT as a consequence.  All numbers below were
measured in these runs; the inputs, outputs and analysis notebooks are kept in
\texttt{combine/Models\_for\_HGKym/} (\texttt{static\_fc10/}, \texttt{diag\_peel\_asym/},
\texttt{diag\_beta899/}).

The RT-scat source used is \texttt{SJChang\_codes/code\_clumpy\_rotating} (clumpy) and
\texttt{code\_smooth\_rotating\_241201} (smooth).  Its run parameters are set by editing
the source before each compilation, so the values found in the source are not
necessarily those used for the paper; only the algorithms are taken from it.

%=============================================================================
\section{Static clumpy sphere: setup}
%=============================================================================

Both codes were run for the same model: a sphere of radius $R_o$ filled with clumps of
radius $r_{\rm cl}=10^{-3}R_o$ at covering factor $f_c=10$ and mean radial column
$N_{\rm HI}=10^{18}\,{\rm cm^{-2}}$ (clump center-to-surface column
$N_{\rm HI,cl}=\tfrac34 N_{\rm HI}/f_c=7.5\times10^{16}\,{\rm cm^{-2}}$), $T=10^4$\,K, no dust,
no bulk motion, the H\,{\sc i} Ly$\alpha$ fine-structure doublet, a central point source
emitting a Gaussian of $\sigma=1\,{\rm km\,s^{-1}}$ around each doublet member, and
$10^6$ photon packets.  Without rotation the model is statistically spherically
symmetric, so every viewing angle must give the same spectrum apart from the scatter of
one clump realization.

Two estimators were used in each code: the image-summed peel-off spectrum of observers at
$\theta=0.1^\circ,15^\circ,\dots,75^\circ,89.9^\circ$ (azimuth $\phi=60^\circ$, the RT-scat
observer set), and escaping photons binned in $\mu=\cos\theta$ over all azimuths
(40 bins in $[-1,1]$).

\begin{center}
\begin{tabular}{lll}
\toprule
 & RT-scat & LaRT\_v2.00 \\
\midrule
clump centers & within $\pm(\Delta x/2-r_{\rm cl})$ of a cell center & uniform over the sphere \\
clumps at $r<0.1R_o$ & none & yes \\
overlap & allowed & not allowed (random sequential adsorption) \\
clump number & $(4/3)f_c(R_o/r_{\rm cl})^2\times(1-0.1^3)$ & $(4/3)f_c(R_o/r_{\rm cl})^2$ \\
peel-off $\tau$ cutoff & 300 & 745.2, now 300 (Section~\ref{sec:lart}) \\
\bottomrule
\end{tabular}
\end{center}

Section~3.2 of the paper states that the clumps are uniformly distributed within the
sphere and do not overlap, which is what LaRT does.

%=============================================================================
\section{Static clumpy sphere: results}
%=============================================================================

\begin{figure}[t]
\centering
\includegraphics[width=\linewidth]{rtscat_static_fc10_angles}
\caption{Spectra at seven viewing angles, normalized to unit area and smoothed with a
25\,km\,s$^{-1}$ boxcar.  Top: LaRT; bottom: RT-scat; left: peel-off; right: $\mu$ bins.}
\label{fig:angles}
\end{figure}

Table~\ref{tab:angles} and Figure~\ref{fig:angles} summarize the angle dependence.
LaRT is isotropic in both estimators, and so are the RT-scat $\mu$ bins.  The RT-scat
peel-off spectra at $\theta=0.1^\circ$ and $89.9^\circ$ are not: they carry 24\% and 12\%
more flux than the mean and a filled-in line center, while $15^\circ$--$75^\circ$ agree with
one another.

\begin{table}[h]
\centering
\caption{Angle dependence of the static $f_c=10$ spectra: integrated flux relative to the
mean over the seven angles, and $\sum|S-\langle S\rangle|/\sum\langle S\rangle$ of the
unit-area spectra.}
\label{tab:angles}
\begin{tabular}{lccc}
\toprule
estimator & flux / mean & shape deviation & peaks [km\,s$^{-1}$] \\
\midrule
LaRT peel-off      & 0.991--1.010 & 0.1--0.6\% & $-43$, $+47$ \\
LaRT $\mu$ bins    & 0.991--1.005 & 0.6--1.5\% & $-43$, $+47$ \\
RT-scat $\mu$ bins & 0.980--1.032 & 1.0--2.6\% & $-43$, $+42$ \\
RT-scat peel-off, $15^\circ$--$75^\circ$ & 0.928--0.932 & 6.4--6.9\% & $-43$, $+47$ \\
RT-scat peel-off, $0.1^\circ$  & 1.238 & 22.8\% & $-33$, $+37$ \\
RT-scat peel-off, $89.9^\circ$ & 1.116 & 10.4\% & $-38$, $+42$ \\
\bottomrule
\end{tabular}
\end{table}

\begin{figure}[t]
\centering
\includegraphics[width=\linewidth]{rtscat_static_fc10_codes}
\caption{LaRT against RT-scat.  Left: all directions ($\mu$ bins summed),
$\sum|L-R|/\sum R=0.027$.  Right: peel-off at $\theta=89.9^\circ$, $\sum|L-R|/\sum R=0.201$.}
\label{fig:codes}
\end{figure}

The angle-averaged spectra of the two codes agree to 2.7\% (Figure~\ref{fig:codes}), but
the $\theta=89.9^\circ$ peel-off spectra differ by 20\%.

\paragraph{Scatterings and run time.}
LaRT scattered $5.92\times10^4$ times per photon, RT-scat $2.85\times10^4$ times
($2.849\times10^4$ and $2.843\times10^4$ for photons born in the two doublet members).
RT-scat ran 45.5\,min on 144 ranks of lart2+lart1 (Xeon Gold 6154); LaRT ran 76.3\,min
with the peel-off cutoff at $\tau=745.2$ and 42.2\,min with it at $\tau=300$, both on 144
ranks of lart4+lart3 (Xeon Gold 6254).  Per scattering, LaRT took 11.1\,$\mu$s (745.2)
and 6.2\,$\mu$s (300) against 13.8\,$\mu$s for RT-scat (on different processors).  The
factor of two in the scattering number is not explained yet; candidates are the empty
center and the clump-free sheets of RT-scat (Section~\ref{sec:sheets}).

%=============================================================================
\section{Problems found in RT-scat}
%=============================================================================

\subsection{Clump-free sheets along the cell faces}
\label{sec:sheets}

RT-scat assigns each clump to a cell of its $201^3$ Cartesian grid and draws its center
within $\pm(\Delta x/2-r_{\rm cl})$ of the cell center (\texttt{set\_clumps\_mod.f90}), so no
clump crosses a cell face.  Thin sheets of width $\sim r_{\rm cl}$ around the planes
$x,y,z={\rm const}$ are therefore nearly free of clumps.  Lines of sight parallel to those
sheets ($\theta\approx0^\circ$: the $x$ and $y$ planes; $\theta\approx90^\circ$: the $z$ planes)
leak more line-center light, which is what the RT-scat peel-off spectra at $0.1^\circ$ and
$89.9^\circ$ show.  The $\mu$ bins average over all azimuths and a range of $\theta$, which
dilutes the effect.

The placement was reproduced in a box of $24^3$ cells (about $4.3\times10^4$ clumps,
\texttt{static\_fc10/rtscat\_clump\_channels.py}); the fraction of rays of length
$0.12R_o$ ($\phi=60^\circ$) that meet no clump is

\begin{center}
\begin{tabular}{lccc}
\toprule
placement & $\theta=0.1^\circ$ & $45^\circ$ & $89.9^\circ$ \\
\midrule
RT-scat (confined to cells) & 0.390 & 0.281 & 0.329 \\
uniform (same number) & 0.301 & 0.296 & 0.293 \\
\bottomrule
\end{tabular}
\end{center}

The clumpy spectra of the paper (Figure~9) are the peel-off spectra of the
$\theta=89.9^\circ$ observer, so they probably contain this leak; this was not checked
against the RT-scat outputs of the paper.

\subsection{Empty center and overlapping clumps}

RT-scat places no clumps at $r<0.1R_o$ and does not prevent overlaps, whereas the paper
(Section~3.2) describes uniformly distributed, non-overlapping clumps.  The clump number
density is that of the whole sphere, so the realized covering factor and mean radial
column are $0.9f_c$ and $0.9N_{\rm HI}$.

\subsection{Smooth sphere: grid sampling of the rotation field}

The smooth rotating sphere is sampled on a $201^3$ Cartesian grid.  Inside $R_i=0.1R_o$
the rotation velocity changes by $\Omega\Delta x=30\,{\rm km\,s^{-1}}$ from one cell to the
next, 2.3 times the thermal width.  Seen from an azimuth that is not a mirror plane of the
grid ($\phi=60^\circ$, used by RT-scat), this staircase makes the edge-on peel-off spectrum
lopsided.  The effect is a property of the grid, not of either code; it was measured with
LaRT (smooth sphere, $N_{\rm HI}=10^{18}\,{\rm cm^{-2}}$, $V_{\rm rot}=300\,{\rm km\,s^{-1}}$,
$\beta=90^\circ$, $10^6$ photons; the AMR row has the fine structure off, the others on):

\begin{center}
\begin{tabular}{lcccc}
\toprule
grid & step in $R_i$ & red/blue flux, $\phi=30^\circ/60^\circ$ & red/blue peak, $\phi=30^\circ/60^\circ$ & $\phi=0^\circ/90^\circ$ flux \\
\midrule
$101^3$ & 59 km\,s$^{-1}$ & 0.936 / 1.071 & 1.054 / 0.958 & 1.007 / 0.998 \\
$201^3$ & 30 km\,s$^{-1}$ & 0.983 / 1.021 & 1.056 / 0.953 & 1.001 / 1.001 \\
$401^3$ & 15 km\,s$^{-1}$ & 1.001 / 1.003 & 1.012 / 0.988 & 1.006 / 0.995 \\
AMR, levels 4--9 & $\le 1.29\,b$ & 0.997 / 0.999 & 0.999 / 0.983 & 1.001 / 0.995 \\
\bottomrule
\end{tabular}
\end{center}

At $\phi=0^\circ$ and $90^\circ$ (mirror planes) the $201^3$ spectrum is symmetric but
rippled (Figure~\ref{fig:grid}); $\phi=30^\circ$ and $60^\circ$ have opposite signs, and
$S_{120}(V)=S_{60}(-V)$ to 1.5\%, as the $x=0$ mirror symmetry of the grid requires.
Switching the fine structure off changes nothing.  The AMR grid
(\texttt{make\_amr\_grid.py --vel\_criterion doppler --vel\_threshold 1.0 --level\_max 9},
$1.1\times10^7$ leaves) reproduces $401^3$ within 0.6--2.3\% with 1/6 of the cells.
The RT-scat smooth curve of the paper used $201^3$ and $\phi=60^\circ$ and is therefore
expected to carry this asymmetry.

\begin{figure}[t]
\centering
\includegraphics[width=\linewidth]{rtscat_diag_azimuth_201}
\includegraphics[width=\linewidth]{rtscat_diag_resolution}
\caption{Smooth rotating sphere.  Top: $201^3$ peel-off spectra at six azimuths (solid)
and their mirror images (dotted).  Bottom: $101^3$, $201^3$, $401^3$ at $\phi=60^\circ$,
$\phi=0^\circ$ and in the $\mu=0$ bin.}
\label{fig:grid}
\end{figure}

\subsection{Direct light seen by a single observer}

The unscattered light of the point source reaches a peel-off observer along one sight
line.  Moving that line by $0.1^\circ$ shifts it at $r=0.5R_o$ by
$1.7\times10^{-3}R_o$, more than a clump radius, so for small $f_c$ the direct spike can
appear in one direction and vanish in the next.  In one LaRT $f_c=1$ realization (same
photons, same clumps), the direct light was 59\% of the flux at $\theta=89.9^\circ,\phi=0^\circ$
and zero at $\theta=90^\circ$; the scattered light differed by less than 1\%.

\subsection{Practical issues}

Run parameters (column density, covering factor, rotation speed, photon number, observer
angles) are set by editing the main program before compiling, so a run cannot be
reproduced from its outputs alone.  The build script compiles \texttt{*\_mod.f90} in
alphabetical order and fails unless the modules happen to be compiled after
\texttt{cons\_mod.f90}; repeating the compile command until all modules are built works.
The observer angles $0.1^\circ$ and $89.9^\circ$ are not needed: the code already treats
$|k_z|=1$ separately and has no singularity at $90^\circ$.

%=============================================================================
\section{Changes made to LaRT}
\label{sec:lart}
%=============================================================================

\begin{itemize}
\item \textbf{Peel-off $\tau$ cutoff.}  The early exit of the photon-transport $\tau$
  integrals used by the peel-off (Cartesian, AMR and clump) is now a single constant,
  \texttt{tau\_huge = 300} in \texttt{define.f90}; it was 745.2 (where $e^{-\tau}$ underflows
  in double precision), defined separately in each routine.  Beyond $\tau=300$,
  $e^{-\tau}<5\times10^{-131}$, so no peel-off weight changes; the static $f_c=10$ run went
  from 76.3 to 42.2\,min with the spectra unchanged (Figure~\ref{fig:taumax}).  Sight-line
  $\tau$ outputs are integrated without the cutoff.
\item \textbf{AMR refinement for velocity fields.}  \texttt{make\_amr\_grid.py} and
  \texttt{AMRGrid} accept \texttt{vel\_criterion='doppler'}: a leaf is refined while the bulk
  velocity changes across it by more than \texttt{vel\_threshold} Doppler widths
  $b=\sqrt{2kT/m+b_{\rm turb}^2}$.
\item \textbf{HDF5 format of the Python writers.}  They write the HDF5\,1.14 format
  (h5py \texttt{libver} upper bound \texttt{'v114'}); with h5py built on HDF5\,2.0,
  \texttt{libver='latest'} produced files that LaRT (linked with HDF5\,1.14) could not read.
\item \textbf{\texttt{par\%grid\_type}} (\texttt{'car'}, \texttt{'amr'}, \texttt{'clump'}) is
  accepted by LaRT\_v2.00 as in LaRT\_v2.10; other values stop the run.
\end{itemize}

\begin{figure}[t]
\centering
\includegraphics[width=\linewidth]{rtscat_static_fc10_taumax}
\caption{Static $f_c=10$ spectra with the peel-off cutoff at $\tau=745.2$ and $\tau=300$
(different random realizations).  Flux ratios 0.993--1.017 at the seven angles, 1.000 over
all directions; shape differences 0.2--1.3\%.}
\label{fig:taumax}
\end{figure}

%=============================================================================
\section{Open questions}
%=============================================================================

\begin{itemize}
\item Why LaRT scatters twice as often as RT-scat in the static $f_c=10$ model.  Running
  LaRT with the RT-scat placement (empty center, overlaps) would separate the
  placement from the transport.
\item Whether the published Figure~9 clumpy curves contain the sheet leak; this needs
  the RT-scat outputs used for the paper.
\end{itemize}

\end{document}
